\documentclass[literature.tex]{subfiles}

\begin{document}
\drama{Čekání na Godota}{Samuel Beckett}
\teorieA{
drama, absurdní drama (tragikomedie o dvou jednáních)
}{
Divadlo absurdity (avantgardní drama 50. let 20. století). Popírá tradiční dramatickou stavbu, dějovost a smysluplný dialog. Blízké existencialismu.
}{
Hovorový, fragmentární jazyk plný opakování, vulgarismů, synonym a slovních her. Dialogy připomínají ping-pong – slova ztrácejí komunikační funkci a slouží jen k zabíjení času. Minimalistické scénické poznámky.
}{
Opakování (anafora, epizeuxis), ironie, černý humor, kontrast mezi banálním a tragickým.
}{
Alegorie (\texthl{Godot = Bůh / naděje / vykoupení}), symbolismus (strom, boty, klobouk, provaz), groteska, oxymóron (tragikomedie).
}{
Dvě jednání, která se téměř zrcadlově opakují. Minimalistická scéna (cesta, strom). Dialogy bez logického postupu, přerušené, nesmyslné. Tragikomický tón – smích z absurdity lidské existence. Absence klasického děje a charakterového vývoje.
}{
VLADIMÍR: Řekni, nechceš se oběsit?\\
ESTRAGON: To bychom mohli. Ale co když se provaz přetrhne?\\
VLADIMÍR: To by byla smůla.\\
ESTRAGON: Ano, to by byla smůla. Ale co když se nepřetrhne?\\
VLADIMÍR: Pak bychom se museli rozejít.\\
ESTRAGON: To by byla smůla.
}
\teorieB{
\wtem{Estragon}{praktický, pesimistický, fyzicky trpící (bolí ho nohy), často zapomíná.}
\wtem{Vladimír}{intelektuálnější, starostlivější, snaží se udržet logiku a naději.}
\wtem{Pozzo}{dominantní, sadistický, bohatý pán; ve druhém jednání slepý a závislý.}
\wtem{Lucky}{ponížený otrok; ve druhém jednání němý. Schopen dlouhého „myšlení“ jen s kloboukem.}
\wtem{chlapec}{posel, který dvakrát oznámí, že Godot nepřijde.}
}{
Dva tuláci Estragon a Vladimír čekají na venkovské cestě u stromu na Godota, který jim má změnit život. Nevědí přesně kde ani kdy. Krátí si čas nesmyslnými dialogy, hádkami, zkoumáním bot a klobouků a úvahami o sebevraždě oběšením. Přijde bohatý Pozzo se svým otrokem Luckym (na provaze). Pozzo předvádí Luckyho „řečnické umění“. Na konci každého jednání přijde chlapec s poselstvím, že Godot dnes nepřijde, ale určitě zítra. Druhé jednání téměř opakuje první – Pozzo je slepý, Lucky němý, strom má pár listů. Tuláci se znovu rozhodnou odejít, ale zůstanou stát.
}{
Dvě jednání (tragikomedie o dvou jednáních). Opakující se struktura (cirkulární kompozice). Prostor: venkovská cesta a strom (minimalistický). Čas: dva po sobě jdoucí dny (neurčitá doba, možná celý život).
}{
Absurdnost lidské existence – marné čekání na smysl, Boha, štěstí nebo vykoupení (\texthl{Godot} jako anagram „God ot“ nebo evokace „God“). Život jako nekonečné zabíjení času v bezvýchodné situaci. Kritika pasivity, rezignace a neschopnosti jednat. Tragikomický obraz lidské bezmoci v poválečné Evropě.
}
\historie{
Napsáno po 2. světové válce (1952 francouzsky, 1953 anglicky). Odráží existenciální krizi, zklamání z války, holokaustu a ztrátu víry v pokrok.
}{}{}{}{
Vrchol divadla absurdity (spolu s Ionescem, Adamovem). Reakce na realismus a tradiční drama. Blízké existencialismu (Camus, Sartre) – absurdity lidského údělu.
}{
(1906 Dublin – 1989 Paříž). Irský spisovatel a dramatik, psal francouzsky i anglicky. Studoval v Dublinu, žil v Paříži. Účastnil se francouzského odboje za 2. světové války. Nobelova cena za literaturu 1969 (převzal ji neochotně). Známý minimalistickým stylem a tématy osamění, paměti a absurdity.
}{
Mezi nejznámější díla patří \textsc{Konec hry}, \textsc{Šťastné dny}, \textsc{Murphy}, \textsc{Molloy} a rozhlasové hry.
}\newpage
\kritika{}{}{
Nadčasové – čekání na „něco lepšího“ (Boha, revoluci, štěstí) zůstává univerzální. Hra rezonuje i v době existenciální úzkosti, sociálních sítí a pasivního konzumu.
}
\nazor{
Čekání na Godota je mistrovské dílo, které na první pohled působí nudně, ale postupně hypnotizuje. Opakující se dialogy a absence děje dokonale vystihují absurdnost života – celý život trávíme čekáním na něco, co nikdy nepřijde. Nejsilnější je tragikomický kontrast: smějeme se banalitám, ale zároveň cítíme mrazení z lidské bezmoci. Godot jako symbol nenaplněné naděje (možná Bůh) je geniální. Beckett ukázal, že drama nemusí mít děj, aby bylo hluboké. Jedno z nejdůležitějších děl 20. století – doporučuji přečíst i zhlédnout.
}
\explain{
\textbl{Divadlo absurdity --} avantgardní směr, který zdůrazňuje nesmyslnost existence, ztrátu smyslu řeči a absence tradičního děje;
\textbl{Tragikomedie --} spojení tragických a komických prvků;
\textbl{Alegorie --} zobrazení abstraktního jevu konkrétním obrazem (Godot = naděje / Bůh).
}
\wpage
\end{document}